% !TeX root = ../papers/nc-p3.tex
\section{非被約三次式を含む非零交換子判別式の全復元}
\label{sec:marked-nonreduced-cubic-commutators}
\begin{theorem}\label{thm:marked-nonreduced-cubic-commutators}
標数零の代数閉体上、三重直線または二重直線と別の単純直線を行列式とする
順序付き跡零三行列 $C$ に対し、全ての $C_Q=C$ は以下で得られる。
各支持直線の法方向の $Q_i$ が階数一行列の跡零部なら圧縮対、零なら一交換子の
二次式で復元する。これらのどちらも起こらない部分は、可逆スライスで割った目標の
残り二作用が既約な場合に限り候補を持つ。候補はただ一つの補助表示になり、
具体的な六座標の線型比較で対称行列 $S$ が得られる。この部分の正確な全逆像は
$S$ の階数3,2,1に応じ二点、空、一つのアフィン二次曲面である。

従って先の各判別式型と合わせ、行列式が恒等零でない全ての順序付き三交換子に
復元手順がある。特に一点を指定した有限窓で $q=3,g>3$ のとき、
正規化した対称側の六行列に対して残る交換子は全判別式型を含め復元できる。
この主張は二切断方程式についての被覆であり、個別の候補の全次数正則性は別に要る。
\end{theorem}
\begin{proof}[補助表示と二つの既約鉛筆]
$\det C_Q\ne0$ なら $\Gamma_Q$ は単射であり、その行消去から得る $A$ は
\emph{重複度も込みで} $\det A$ が $\det C_Q$ の定数倍になる。
実際 $\det A(v)$ は、九空間において $\Gamma_Q$ の六列と $v\otimes U$ の三列を並べた
行列式の定数倍である。$v=s\wedge t$ のときこれは $s,t$ への評価で
$D^*\otimes U$ を $(v\otimes U)$ で割った六次元商への写像の行列式となる。
体積による同定で商の行列式の倍率は一定である
（$s,t$ を基底の最初の二元に拡張し外積の基底変換を取ればよい）。
定理\ref{thm:rank-eight-finite-symmetry}の六変数恒等式からこれは $\det[Q_s,Q_t]$ である。

座標を $\det A$ が $x^3$ または $x^2(x+y)$ の倍になるように取り、
$A=xI+yJ+zD$ と規格化する。二作用が可約なら同時三角化できる。実際不変直線か不変平面の
残る二次元部分でも全行列式は一次式の積である。その線型固有文字の一つを引くと
二スライスの包は階数1以下で、共通核か共通像により同時不変直線がある。
従ってもとの鉛筆に共通左固有ベクトルがあり、行空間中の階数一テンソルを与える。
補題\ref{lem:marked-squarefree-symmetrizer}末尾の階数一テンソルと $Q_n$ の同値の証明は
被約性を用いず、ここでは $Q_n=0$ も許す。逆にこのケースでは必ず可約となる。

既約な場合の左右基底同値類は、固定した $x,y,z$ のままで次の各一変数 $k\ne0$ で尽きる。
\[
\begin{array}{c|c|c}
 &J&D\\ \hline
x^3&E_{12}+E_{23}& k(E_{21}-E_{32})\\
x^2(x+y)&E_{12}+E_{33}&
\left(\begin{smallmatrix}k&-k&0\\0&-k&1\\k^2&-k^2&0\end{smallmatrix}\right)
\end{array}
\]
この表の行列単位だけは $1,2,3$ 始まりである。
三重直線では $J,D$ の全結合は冪零であり、一般元は階数2である（階数1以下の空間は可約）。
まずこの一般元を $J=E_{12}+E_{23}$ として計算する。
$\operatorname{tr}D=\operatorname{tr}(JD)=\operatorname{tr}(J^2D)=0$ から
$D=\left(\begin{smallmatrix}a&b&c\\d&e&f\\0&-d&-a-e\end{smallmatrix}\right)$。
既約性で $d\ne0$、$\operatorname{tr}(JD^2)=d(2a+e)=0$。
$I+sJ$ の共役で対角を消すと $\operatorname{tr}D^2=0$ から $b=f$、
$I+tJ^2$ の共役で両方消え、$\det D=0$ から $c=0$ となる。
得た型ではどの非零の結合も階数2なので、最初から指定した $y$ 方向が使える。
$k$ は $[J,D]$ の二重固有値で一意である。

二重直線でも最初は一般の $y$ 方向（$y$ の係数を1にした $J+tD$）を使う。
二重零固有値が半単純なら階数1、一般にそうなら上と同じく可約である。
従って $J=E_{12}+E_{33}$ にできる。三つの跡条件から
$D=\left(\begin{smallmatrix}a&b&u\\0&-a&v\\w&h&0\end{smallmatrix}\right)$。
$\operatorname{tr}D^2=\operatorname{tr}(JD^2)=0$ から
$uw+vh=-a^2$、$vw=a^2$ を得る。$v=0$ なら第二座標行が共通左固有、
$w=0$ なら第一座標列が共通右固有なので $a,v,w\ne0$。
$J$ の中心化群の上三角剪断とブロック間の倍率で $u=0,v=1$ とし、
$\det D=0$ を用いると上表の $k=a$ となる。残余中心化群は $a$ を変えない。
この型の $J+tD$ は全 $t$ で階数2であり、指定した $y$ に戻せる。
二つの表の型が既約であることも $J$ の左右固有直線を $D$ が保たないことから分かる。
\end{proof}

\begin{proof}[固定基底での逆像]
二つの既約鉛筆の六空間性は定理\ref{thm:marked-elliptic-commutator-doubling}の
双対三テンソルの証明に従う。そこで必要なのは二作用の既約性だけである。
参照目標を $M=xI+yJ'+zD'$ とすると、二型の順に
\[
\begin{split}
 J'&=E_{12}-E_{23}/2,\quad D'=-k E_{21}/2-k E_{32},\\
 J'&=\begin{pmatrix}0&1&-1/k\\0&1&-1/(2k)\\0&0&0\end{pmatrix},\quad
 D'=\begin{pmatrix}0&-k&1\\ k/2&-k&1/2\\0&-k^2&k\end{pmatrix}.
\end{split}
\]
この表示の抽象パラメータは二型とも $-k/2$。
第一の場合は $\operatorname{diag}(1,1,-2)$、
第二では列 $(1,0,0)^t,(0,-1,-2k)^t,(1,1,0)^t$ の基底で $J'$ を表の形に戻すと分かる。
よって既約な目標を規格化すれば $k$ は唯一に求まる。
前節と同じ六座標 $(P,T,U,A,B,D_0)$ と自由列の順で六空間に以下を代入する。
\[
\begin{split}
 &(-U,kP+D_0,B,kA,D_0,kT),\\
 &(-2k^2P+k^2A-kD_0+U,-kP+kT-B+D_0,P-A,-kA+kT-D_0,A-T,k^2T).
\end{split}
\]
これらは $k\ne0$ で同型である。$C_Q(e_0)$ の六余因子係数を $\Psi$ とすると
$C_Q(v)=\Psi(\operatorname{adj}Z)M(v)$ で、六行 $0,1,2,3,4,6$ の小行列式は
$4k^{10},-4k^{10}$ である。\path{computations/verify_nonreduced_commutators.py} で全係数を検算する。
$C=L M R$ の基底比較の自由度は逆倍率のみである。実際 $M(e_0)=I$ なので、
二つの比較の差は $J',D'$ の共通中心化に属し、既約性によりスカラーとなる。
従って $\operatorname{adj}Z=\Psi^{-1}(RL)$ による実際の解とその階数分類が従う。

補正がある各支持直線で法方向を最初の添字に取れば、
$C_1K+C_2K\subset kw$ による全ての圧縮対と定理\ref{thm:marked-transverse-conic-commutators}の
線型復元が使える。直線上の一般階数が1以下でも、全 $K\in\mathbf P(U^*)$ について
$C_1K+C_2K$ の次元を小行列式で判定すればよい。次元1では $kw$ はその像、
次元0では全 $[w]\in\mathbf P(U)$、次元2以上なら解なしである。
$Q_0=0$ なら $C_1=C_2=0$ で三重直線型に限られる。
逆にその入力では $C_0$ が可逆である。$Q_0$ が非スカラーなら、
巡回の場合その中心化環は可換、非巡回の場合その中心化環は
階数一補正の像直線を保つため、いずれもそこにある二元の交換子が可逆になることはない。
ゆえに $Q_0=0$ が必要である。残りは巡回な跡零 $X$ について
$\operatorname{tr}(XC_0)=\operatorname{tr}(X^2 C_0)=0$ の全てを取り、
$Q_1=X$、$[X,Q_2]=C_0$ の跡零解のアフィン二平面を取れば尽くされる。
最後に有限窓では定理\ref{thm:no-three-bracket-compression}とその直前の可換交叉の排除、
および定理\ref{thm:marked-two-commutators}が $q=3,g>3$ の非零判別式または二交換子復元を保証する。
\end{proof}
